\documentclass[conference]{IEEEtran}

\usepackage{cite}
\usepackage{amsmath}
\usepackage{graphicx}

\begin{document}

%================ TITLE =================%
\title{\textbf{ROUTEMATE}}

%========= SUBTITLE (LIKE FIRST IMAGE) =========%
\makeatletter
\def\@maketitle{
\newpage
\null
\vskip 1em%
\begin{center}%
{\LARGE \bfseries ROUTEMATE \par}%
\vskip 0.5em%
{\large \textit{An Innovative Ride-Sharing System for Sustainable Urban Mobility}\par}%
\vskip 1.5em%

%================ AUTHORS ROW (LIKE FIRST IMAGE) =================%
\begin{tabular}{ccccc}

\begin{tabular}{c}
\textit{Ms. Anitta K Mathew} \\
Assistant Professor \\
Department of IT \\
VJCET
\end{tabular}

&

\begin{tabular}{c}
\textit{Anupam Prakash} \\
UG Scholar \\
Department of IT \\
VJCET
\end{tabular}

&

\begin{tabular}{c}
\textit{Georgit Dain} \\
UG Scholar \\
Department of IT \\
VJCET
\end{tabular}

&

\begin{tabular}{c}
\textit{Jeswin Jaison} \\
UG Scholar \\
Department of IT \\
VJCET
\end{tabular}

&

\begin{tabular}{c}
\textit{Johns Joseph} \\
UG Scholar \\
Department of IT \\
VJCET
\end{tabular}

\end{tabular}

\end{center}
\par
\vskip 1.5em}
\makeatother

\maketitle

%================ ABSTRACT =================%
\begin{abstract}
This paper presents RouteMate, a next-generation ride-sharing platform designed to address urban traffic congestion and excessive fuel consumption through optimized vehicle usage. The system enables private vehicle owners to share rides with passengers whose destinations align with their intended routes using real-time, route-based matching. Unlike traditional ride-sharing platforms that rely on fixed schedules or commercial fleets, RouteMate facilitates spontaneous and flexible ride coordination for daily commuters. The platform integrates Aadhaar-based user verification to ensure security, trust, and accountability among users. Additionally, advanced location tracking and dynamic route optimization algorithms are employed to minimize detours, reduce travel time, and improve ride efficiency. To encourage long-term participation, RouteMate incorporates an incentive-driven reward mechanism that provides fuel vouchers, carbon credit points, and user discounts. By maximizing vehicle occupancy and reducing single-passenger travel, the system contributes to sustainable transportation, reduced emissions, and a smarter urban mobility ecosystem.
\end{abstract}

%================ INDEX TERMS =================%
\begin{IEEEkeywords}
RouteMate, Ride-Sharing, Urban Mobility, Dynamic Route Optimization, Aadhaar Verification, Sustainable Transportation, Smart Mobility
\end{IEEEkeywords}

%================ INTRODUCTION =================%
\section{INTRODUCTION}

Urban transportation systems are facing unprecedented challenges due to rapid population growth, increased private vehicle ownership, and inefficient utilization of road infrastructure. Traffic congestion, fuel overconsumption, rising carbon emissions, and prolonged commuting times have become defining characteristics of modern metropolitan environments. A significant portion of these issues stems from underutilized private vehicles, where a majority of cars operate with vacant seats, contributing unnecessarily to congestion and environmental degradation.

Existing ride-sharing platforms such as Uber and BlaBlaCar have introduced shared mobility models to mitigate these problems. However, many of these systems primarily rely on commercial driver fleets, scheduled bookings, or static ride matching mechanisms. Such approaches often fail to dynamically optimize real-time route alignment between private vehicle owners and passengers traveling in similar directions. As a result, detours, inefficiencies, and scalability limitations persist.

Furthermore, traditional ride-sharing applications emphasize convenience and cost reduction but rarely integrate sustainability-driven mechanisms such as carbon footprint tracking, eco-incentive systems, or intelligent detour minimization algorithms. The absence of a fully dynamic, route-aligned, and incentive-driven ride-sharing model creates a gap in achieving truly sustainable urban mobility.

Recent advancements in GPS tracking, real-time mapping APIs, geospatial analytics, and intelligent matching algorithms have enabled the development of adaptive transportation systems capable of dynamic route optimization. Techniques such as shortest-path computation, geofencing, proximity analysis, and deviation threshold modeling allow efficient identification of compatible ride matches without significantly altering a driver’s original route.

This paper presents both a conceptual study of existing ride-sharing limitations and a proposal for an enhanced system named RouteMate. The proposed system employs a Multi-Layer Smart Mobility Framework comprising:
\begin{itemize}
\item Real-Time Route Monitoring
\item Intelligent Route Deviation Analysis
\item Secure User Verification
\item Eco-Incentive and Reward Mechanism
\end{itemize}

This layered architecture enables dynamic ride matching between private drivers and passengers while ensuring minimal route deviation, enhanced safety, and sustainability-driven participation. Unlike conventional models, RouteMate emphasizes spontaneous ride sharing based purely on route alignment rather than scheduled pooling.

The motivation behind this work arises from key challenges identified in existing systems — inefficient seat utilization, lack of real-time adaptive route matching, minimal integration of sustainability incentives, and limited community-driven mobility frameworks. By addressing these gaps, RouteMate aims to establish a scalable, environmentally conscious, and intelligent ride-sharing ecosystem that supports smart city initiatives.

The subsequent sections of this paper provide a comprehensive review of related works, analysis of their strengths and limitations, detailed system architecture, algorithmic design, and evaluation of the proposed framework.

%================ RELATED WORK (ONLY HEADING AS REQUESTED) =================%
\section{RELATED WORK}

Ride-sharing and urban mobility optimization have been extensively studied, with existing solutions focusing on commercial fleet management, fixed carpool scheduling, and static ride-matching algorithms. While these systems improve transportation efficiency, they often lack real-time dynamic route alignment and sustainability-driven incentive mechanisms. The following works highlight key contributions and limitations in this domain.

\textbf{A. Route-Based Ride Sharing in Uber Technologies Inc.}

\vspace{2pt}
This work enables shared mobility by matching multiple passengers traveling along similar routes using real-time GPS tracking and dynamic route optimization algorithms. The system clusters ride requests based on proximity and destination alignment, reducing travel cost and traffic congestion in urban areas. It effectively utilizes live traffic data and shortest-path algorithms to assign passengers to available drivers. However, this approach is primarily designed for commercial fleet operations and profit-oriented ride allocation. It does not prioritize minimal route deviation for private vehicle owners and lacks sustainability-driven reward mechanisms for reducing carbon emissions. In contrast, the proposed RouteMate system focuses on real-time route-based matching for private vehicle users by calculating deviation thresholds and additional travel time before confirming a ride. It also integrates an eco-incentive model that encourages fuel savings and carbon footprint reduction, making it more suitable for community-driven sustainable urban mobility.

\vspace{4pt}
\textbf{B. Long-Distance Carpooling Platforms: A Survey}

\vspace{2pt}
This survey highlights existing carpooling systems that facilitate ride sharing through pre-scheduled trip postings and user-based booking mechanisms, focusing on trust models, rating systems, and cost-sharing strategies. These platforms allow drivers to publish planned journeys while passengers reserve available seats, improving vehicle occupancy and reducing travel expenses. The approach is particularly effective for intercity and long-distance travel where schedules are predetermined. However, it relies heavily on advance trip planning and does not support spontaneous, real-time intra-city ride matching. Additionally, route optimization is typically static and does not dynamically evaluate minimal deviation thresholds during live travel.

In contrast, the proposed RouteMate system addresses these limitations by enabling real-time, GPS-based dynamic route alignment for urban commuting. Instead of relying on pre-scheduled journeys, RouteMate continuously analyzes live driver routes and calculates acceptable deviation thresholds before confirming matches. This approach ensures minimal detour, faster ride confirmation, and greater suitability for everyday city travel while promoting sustainable shared mobility.

\vspace{4pt}
\textbf{C. AI-Driven Ride Matching and Optimization in Ola Cabs}

\vspace{2pt}
This work introduces an AI-based ride-sharing framework that enhances passenger clustering and route optimization using machine learning techniques within the platform’s operational architecture. The system analyzes user demand patterns, traffic conditions, and route similarities to allocate shared rides efficiently, improving travel time estimation and fleet utilization. While this approach is effective for real-time ride allocation and congestion reduction, it primarily supports commercial fleet operations and centralized pricing models. Additionally, the effectiveness of the system depends on large-scale user data and continuous model training to maintain high matching accuracy.

In contrast, the proposed RouteMate system focuses on peer-to-peer, route-based ride sharing for private vehicle owners by dynamically calculating route deviation thresholds and additional travel time before confirming matches. Rather than relying solely on large-scale predictive models, RouteMate emphasizes real-time route alignment and sustainability incentives such as fuel savings and carbon reduction benefits. This makes the system more suitable for community-driven urban mobility while addressing limitations associated with centralized fleet-based AI models.

\vspace{4pt}
\textbf{D. Route Optimization and Traffic-Aware Navigation in Google Maps Platform}

\vspace{2pt}
This work presents a route optimization approach based on real-time traffic monitoring and shortest-path algorithms, enabling efficient navigation through congestion-aware routing. The system utilizes algorithms such as Dijkstra’s and A* to compute optimal travel paths while continuously analyzing live traffic data to update route recommendations. Although this method achieves high accuracy in navigation and travel time estimation, it is primarily designed for individual route guidance rather than peer-to-peer ride matching. Additionally, it does not evaluate ride feasibility based on route deviation thresholds or shared mobility optimization.

In contrast, the proposed RouteMate system extends traditional route optimization by integrating a dynamic ride-matching engine that analyzes route overlap, calculates acceptable deviation limits, and estimates additional travel time before confirming a match. By combining traffic-aware navigation with real-time peer matching and sustainability incentives, RouteMate addresses the limitations of standalone navigation systems and supports efficient, community-driven ride sharing.

\vspace{4pt}
\textbf{E. Real-Time Traffic Data Integration in Google Maps and Shared Mobility Platforms}

\vspace{2pt}
This work proposes a traffic-aware mobility approach by analyzing real-time road conditions, vehicle density, and route congestion patterns using continuous GPS data streams and traffic analytics. The system monitors travel speed variations, road blockages, and route efficiency metrics to dynamically adjust navigation paths and improve travel time accuracy. While effective for optimizing individual travel routes and reducing congestion impact, this approach relies heavily on continuous real-time data processing and does not inherently support collaborative ride matching after route computation. It also lacks mechanisms for evaluating seat availability, route deviation feasibility, or sustainability-based incentives within shared mobility contexts.

In contrast, the proposed RouteMate system extends real-time traffic analysis by integrating a route-based peer matching mechanism that evaluates route overlap, acceptable detour thresholds, and additional travel time before confirming ride sharing. By combining traffic-aware optimization with sustainability incentives such as fuel savings and carbon footprint reduction, RouteMate provides a comprehensive, community-driven solution that enhances both mobility efficiency and environmental responsibility beyond traditional navigation systems.

\vspace{4pt}
\begin{table}[!h]
\centering
\footnotesize
\renewcommand{\arraystretch}{1.3}
\setlength{\tabcolsep}{2pt}

\begin{tabular}{|c|p{1.8cm}|p{2.6cm}|p{2.6cm}|}
\hline
\textbf{S.No} & \textbf{Title} & \textbf{Advantages} & \textbf{Disadvantages} \\
\hline


1 &
Route-Based Ride Sharing in Uber Technologies Inc. &
Uses real-time GPS tracking and dynamic clustering of passengers based on route alignment; reduces congestion and travel cost. &
Designed mainly for commercial fleet operations; lacks minimal deviation logic for private vehicles and sustainability incentives. \\
\hline

2 &
Long-Distance Carpooling Platforms: A Survey &
Improves vehicle occupancy, trust models, rating systems, and cost sharing; effective for scheduled and intercity travel. &
Relies on pre-scheduled trips; lacks real-time intra-city matching and dynamic deviation analysis. \\
\hline

3 &
AI-Driven Ride Matching and Optimization in Ola Cabs &
Uses machine learning for ride clustering, route optimization, and travel time estimation; improves fleet utilization. &
Depends on large-scale data and centralized fleet models; not suitable for peer-to-peer private ride sharing. \\
\hline

4 &
Route Optimization and Traffic-Aware Navigation in Google Maps Platform &
Accurate shortest-path computation using Dijkstra’s and A* with real-time traffic-aware navigation. &
Focused on individual navigation; lacks ride-matching feasibility and deviation threshold evaluation. \\
\hline

5 &
Real-Time Traffic Data Integration in Google Maps and Shared Mobility Platforms &
Analyzes traffic density and live GPS patterns to improve route efficiency and travel time accuracy. &
Requires continuous real-time processing; lacks collaborative ride matching and sustainability incentive mechanisms. \\
\hline

\end{tabular}

\vspace{2pt}
\caption{Significant Insights from the Literature}
\end{table}

\section{GAPS IDENTIFIED FROM THE LITERATURE}

\textbf{1) Dependence on Commercial Fleet Models –} Most existing ride-sharing systems are designed primarily for commercial taxi or fleet-based operations, limiting participation from private vehicle owners and reducing community-driven mobility opportunities.

\vspace{4pt}

\textbf{2) Lack of Real-Time Route Deviation Analysis –} Many platforms rely on pre-scheduled trips or basic destination similarity without dynamically calculating acceptable route deviation thresholds and additional travel time, which can discourage driver participation.

\vspace{4pt}

\textbf{3) Limited Support for Spontaneous Intra-City Travel –} Several existing carpooling models focus on long-distance or pre-planned journeys and do not effectively support real-time, on-demand ride matching within urban environments.

\vspace{4pt}

\textbf{4) Absence of Sustainability-Centered Incentives –} Most current solutions emphasize cost-sharing and operational efficiency but do not incorporate structured reward mechanisms such as fuel-saving benefits or carbon emission reduction incentives to encourage environmentally responsible travel behavior.

\section{PROPOSED METHODOLOGY}

\textf{A. SYSTEM ARCHITECTURE}

Based on the detailed analysis of existing ride-sharing systems and the research gaps identified in the literature survey, this paper proposes RouteMate, a real-time route-based ride-sharing platform designed to optimize private vehicle utilization and promote sustainable urban mobility. The framework aims to overcome key limitations of previous systems such as dependence on commercial fleets, lack of dynamic route deviation analysis, and absence of sustainability-driven incentives.

RouteMate follows a layered modular architecture consisting of five layers: a Presentation Layer built using a Flutter-based mobile application for user registration, ride requests, live tracking, and reward visualization; an Application Logic Layer implementing the route-matching engine, deviation threshold calculation, and travel time estimation algorithms; a Location Monitoring Layer responsible for real-time GPS tracking and route alignment using map APIs; a Communication Layer enabling secure API interactions between the mobile application and backend server using HTTPS protocols; and a Data Persistence Layer utilizing a cloud-based database such as MongoDB or MySQL for storing user profiles, ride history, route logs, and reward records. This structured architecture ensures scalability, secure communication, efficient ride matching, and seamless integration of sustainability incentives within the RouteMate ecosystem.

\textbf{Key Components of the Framework}

\textbf{1) Real-Time Route Monitor (Layer 1 – Live Location Tracking):} The Route Monitor module utilizes GPS services and map APIs to continuously track the driver’s live location and predefined travel route. A dynamic route-alignment algorithm evaluates nearby ride requests and triggers a Feasible Match Alert when a rider’s pickup and drop locations fall within a predefined deviation threshold (e.g., 2 km detour or 5 minutes additional travel time). This ensures that ride matching does not significantly alter the driver’s intended path.

\vspace{4pt}

\textbf{2) Route Deviation Engine (Layer 2 – Mathematical Optimization):} The system applies shortest-path algorithms such as Dijkstra’s or A* to compute optimal routes and estimate additional travel distance and time. The deviation cost function is calculated as:
\[
\text{Deviation} = (\text{NewRouteDistance} - \text{OriginalRouteDistance})
\]
If the deviation remains within the acceptable threshold, the ride request is marked as suitable. This mathematical validation ensures efficient and fair ride allocation while minimizing inconvenience to drivers.

\vspace{4pt}

\textbf{3) User Identity Trust Verification (Layer 3 – Secure Authentication):} All users undergo secure authentication using OTP verification and profile validation. Driver and passenger details are cross-verified, and ratings are maintained to ensure safety and trust. Suspicious accounts or inconsistent ride behaviors trigger administrative review mechanisms, enhancing platform reliability.

\vspace{4pt}

\textbf{4) Sustainability Incentive Engine (Layer 4 – Eco-Scoring Mechanism):} A weighted eco-score model assigns reward points based on shared distance, number of co-passengers, and estimated carbon emission savings. For example:
Shared Distance (4 pts), Number of Riders (3 pts), and Reduced CO Emission Estimate (3 pts). When cumulative eco-points exceed a defined threshold, users become eligible for fuel vouchers or carbon credit incentives. This multi-factor scoring approach encourages sustainable ride-sharing participation.

\vspace{6pt}

\textbf{Advantages of the Proposed Framework}

\begin{itemize}
\item Enables real-time, deviation-aware ride matching, ensuring minimal disruption to the driver’s planned route.
\item Supports spontaneous intra-city ride sharing using continuous GPS tracking and dynamic route evaluation.
\item Encourages environmentally responsible behavior through a structured eco-reward and carbon-saving incentive system.
\item Enhances user safety through secure authentication, rating systems, and monitoring mechanisms.
\item Reduces traffic congestion and fuel consumption by improving private vehicle seat utilization.
\item Follows a modular layered architecture that improves scalability, maintainability, and future integration of AI-based predictive matching.
\item Ensures smooth application performance through optimized API communication and efficient backend processing, preventing delays during real-time matching operations.
\end{itemize}

%================ SPACE FOR Architecture Diagram of the Proposed System =================%
\vspace{3pt}
\textf{B. ALGORITHM}

\textbf{Route-Based Dynamic Ride Matching Algorithm}

The proposed system performs real-time ride matching through route deviation analysis, shortest-path optimization, trust validation, and sustainability scoring. The step-by-step procedure is described as follows.

\textbf{Step 1 (Start):} Initialize the system, authenticate the user (driver/passenger), and enable GPS location services. The driver enters the destination and confirms route generation.

\textbf{Step 2:} The system retrieves the optimal original route using shortest-path algorithms (Dijkstra/A*) via the map API and stores the baseline distance and estimated travel time.

\textbf{Step 3:} The Live Route Monitor activates continuous GPS tracking and updates the driver’s current position at regular intervals.

\textbf{Step 4:} When a passenger submits a ride request, the system captures pickup location, drop location, and timestamp.

\textbf{Step 5:} The Route Deviation Engine computes a new temporary route incorporating the passenger’s pickup and drop points. The deviation is calculated as:
\[
\text{Deviation} = (\text{NewRouteDistance} - \text{OriginalRouteDistance})
\]

\textbf{Step 6:} The additional travel time is estimated. If the deviation is within the predefined ThresholdDistance and the additional travel time is within ThresholdTime, the ride request is marked as Feasible.

\textbf{Step 7:} The Trust Verification module validates user authentication status, profile rating, and ride history to ensure safety compliance.

\textbf{Step 8:} The Sustainability Engine estimates shared distance and calculates eco-points based on reduced fuel usage and CO emission savings.

\textbf{Step 9:} The ride request is classified into three categories: 0 – Not Feasible, 1 – Feasible, and 2 – High Priority Match (minimal deviation with high eco-benefit).

\textbf{Step 10:} If classified as Feasible or High Priority, a confirmation notification is sent to both driver and passenger through secure API communication.

\textbf{Step 11:} Upon acceptance, the system updates ride records in the database, including route data, deviation value, travel time, and eco-points earned.

\textbf{Step 12:} The dashboard updates ride statistics, fuel savings estimation, and carbon reduction metrics for user analytics and monitoring.

\textbf{Step 13 (Stop):} The system continues monitoring until the ride is completed or the driver ends the session.

\section{RESULTS}

The proposed RouteMate system was evaluated through functional testing, simulated ride scenarios, and real-time route matching analysis to assess its effectiveness in improving urban mobility and vehicle utilization. The system successfully demonstrated real-time route-based ride matching between drivers and passengers traveling in similar directions with minimal deviation from the original route.

Experimental observations indicate that the dynamic route alignment mechanism efficiently identifies feasible ride matches by computing deviation distance and additional travel time using shortest-path algorithms. In most simulated intra-city travel scenarios, the deviation remained within acceptable thresholds (2 km or 5 minutes), ensuring that drivers’ planned routes were not significantly disrupted while still enabling effective ride sharing.

The Live Route Monitoring module, integrated with GPS tracking and map APIs, provided continuous location updates and accurate route alignment analysis. This enabled spontaneous ride matching without requiring pre-scheduled trips, making the system highly suitable for daily urban commuting. Additionally, the Trust Verification module ensured secure user participation through authentication, profile validation, and ride history tracking, thereby enhancing platform safety and reliability.

The Sustainability Incentive Engine produced measurable eco-benefits by estimating shared travel distance and carbon emission reduction. Users accumulated eco-points based on ride-sharing activity, which can be translated into fuel vouchers or carbon credit incentives. This mechanism significantly improved user engagement and encouraged environmentally responsible travel behavior.

Overall, the results confirm that RouteMate effectively reduces single-occupancy vehicle usage, improves seat utilization, minimizes unnecessary detours, and promotes sustainable shared mobility in urban environments.

\section{DISCUSSIONS}

The implementation of RouteMate highlights the importance of real-time, route-aligned ride sharing as a scalable solution for modern urban transportation challenges. Unlike conventional ride-hailing platforms that rely on commercial fleets and centralized allocation, RouteMate focuses on peer-to-peer ride matching among private vehicle users, thereby expanding participation and improving community-driven mobility.

One of the key strengths of the system lies in its deviation-aware matching strategy, which ensures that ride feasibility is determined based on route overlap, additional travel time, and acceptable detour limits. This approach addresses a major limitation in traditional carpooling systems that rely primarily on static destination similarity or pre-scheduled journeys.

Furthermore, the integration of sustainability-centered incentives distinguishes RouteMate from existing platforms. By incorporating eco-scoring based on shared distance and carbon emission savings, the system promotes environmentally conscious commuting while simultaneously reducing traffic congestion and fuel consumption. This aligns with smart city initiatives and sustainable urban development goals.

However, certain limitations must be considered. The accuracy of real-time ride matching depends on continuous GPS tracking, map API responsiveness, and network connectivity. In regions with poor internet availability or GPS inaccuracies, route alignment and feasibility evaluation may experience minor delays. Additionally, large-scale deployment would require robust backend infrastructure to handle high volumes of ride requests and real-time location data processing.

Despite these limitations, the proposed layered architecture ensures scalability, modular integration, and future extensibility. The framework can be further enhanced by incorporating AI-based predictive matching, demand forecasting, and adaptive incentive optimization to improve matching accuracy and user participation. Therefore, RouteMate presents a practical and sustainable approach to intelligent urban ride sharing with strong potential for real-world implementation.

\section{CONCLUSION}

This paper presented RouteMate, a real-time route-based ride-sharing system, along with a comprehensive survey of existing shared mobility platforms and route optimization approaches. The survey examined commercial fleet-based ride-sharing models, pre-scheduled carpooling systems, and traffic-aware navigation solutions, highlighting their strengths and limitations. From this analysis, it was observed that while existing systems effectively reduce travel costs and improve route efficiency, they largely depend on commercial operations, lack dynamic deviation-aware peer matching, and provide limited sustainability-driven incentives.

The proposed RouteMate framework addresses these limitations by integrating a layered architecture comprising Real-Time Route Monitoring, Mathematical Route Deviation Analysis, Secure User Verification, and a Sustainability-Based Incentive Engine. This combination enables efficient peer-to-peer ride matching with minimal detours, promotes optimal seat utilization in private vehicles, and encourages environmentally responsible travel through structured eco-reward mechanisms.

Furthermore, the system’s modular and scalable design ensures seamless communication between the mobile application and backend services, enabling smooth real-time matching and analytics updates. Unlike conventional ride-sharing platforms that prioritize fleet efficiency, RouteMate introduces a community-driven, sustainability-focused mobility solution suitable for spontaneous intra-city travel.

Future work may focus on integrating AI-based predictive demand analysis, blockchain-based carbon credit tracking, support for electric vehicle prioritization, and large-scale deployment within smart city infrastructures to further enhance urban transportation efficiency and environmental impact reduction.

\begin{thebibliography}{00}

\bibitem{b1} \textit{H. Yu, X. Jia, H. Zhang, and J. Shu}, “Efficient and privacy-preserving ride matching using exact road distance in online ride-hailing services,” \textit{IEEE Transactions on Services Computing}, vol. 15, no. 4, pp. 1841–1854, 2022, doi: 10.1109/TSC.2020.3022875.

\bibitem{b2} \textit{M. B. Amasyali and E. Gul}, “VoIP integration for mobile ride-sharing application,” in \textit{Proc. 7th Int. Conf. Communication Systems and Network Technologies (CSNT)}, 2017, pp. 44–47, doi: 10.1109/CSNT.2017.8418509.

\bibitem{b3} \textit{M. T. Brillian, Y. G. Sucahyo, Y. Ruldeviyani, and A. Gandhi}, “Revealing the misuse of motorcycle ride-sharing applications using extended deterrence theory approach,” in \textit{Proc. Int. Conf. Advanced Computer Science and Information Systems (ICACSIS)}, 2018, pp. 131–135, doi: 10.1109/ICACSIS.2018.8618258.

\bibitem{b4} \textit{S. Ma, Y. Zheng, and O. Wolfson}, “Real-time city-scale taxi ridesharing,” \textit{IEEE Transactions on Knowledge and Data Engineering}, vol. 27, no. 7, pp. 1782–1795, 2015, doi: 10.1109/TKDE.2014.2334313.

\bibitem{b5} \textit{Z. Liu, Z. Gong, J. Li, and K. Wu}, “mT-Share: A mobility-aware dynamic taxi ridesharing system,” \textit{IEEE Internet of Things Journal}, vol. 9, no. 1, pp. 182–198, 2022, doi: 10.1109/JIOT.2021.3102638.

\bibitem{b6} \textit{J.-P. Li, G.-J. Horng, S.-T. Cheng, and C.-F. Chen}, “Intelligent ridesharing system for taxi to reduce cab fee,” in \textit{Proc. IEEE 12th Int. Conf. Networking, Sensing and Control}, 2015, pp. 468–473, doi: 10.1109/ICNSC.2015.7116082.

\bibitem{b7} \textit{M. I. Makarim, N. Selviandro, and G. S. Wulandari}, “Route recommendation simulation for ride-sharing autonomous vehicle: A comparative study of A* and Dijkstra algorithm,” in \textit{Proc. 1st Int. Conf. Software Engineering and Information Technology (ICoSEIT)}, 2022, pp. 216–221, doi: 10.1109/ICoSEIT55604.2022.10029975.

\bibitem{b8} \textit{K. Zhou}, “The formation mechanism and characteristics of urban multimodal traffic system based on the interaction between urban spatial mode and traffic development,” in \textit{Proc. Int. Conf. Urban Engineering and Management Science (ICUEMS)}, 2020, pp. 24–28, doi: 10.1109/ICUEMS50872.2020.00014.

\end{thebibliography}

\end{document}